% Canonical Paper I appendix.
This appendix records sign checks and coordinate calculations used in
Sections~8 and~9.  The chronological convention is unchanged: entries in a
history are performed from left to right.  The ACS is oriented by
\(dA\wedge dM\), and the contact state space is oriented by
\(du\wedge da\wedge dv\).

\subsection{The oriented ACS rectangle}

Let
\[
    \gamma=(\mathsf A_p,\mathsf M_q),
    \qquad
    \delta=(\mathsf M_q,\mathsf A_p).
\]
Their common terminal multiplicative charge is \(M_*=q\).  The chain
\(C_\gamma-C_\delta\) traverses the bottom edge from \((0,0)\) to \((p,0)\),
the right edge to \((p,q)\), the top edge back to \((0,q)\), and the left edge
back to the origin.  For \(p,q>0\) this is the positive boundary convention
for \(dA\wedge dM\).  Since vertical edges do not contribute to
\(\eta_*=e^{q-M}dA\),
\begin{align*}
    \int_{C_\gamma-C_\delta}\eta_*
      &=\int_0^p e^q\,dA+\int_p^0 1\,dA\\
      &=p(e^q-1).
\end{align*}
On the other hand,
\[
    \int_{[0,p]\times[0,q]}e^{q-M}\,dA\wedge dM
      =p\int_0^q e^{q-M}\,dM
      =p(e^q-1).
\]
This verifies both the sign and the target-frame weight.

The same formula holds for signed \(p\) and \(q\).  More precisely, regard the
rectangle as the oriented singular chain parametrized by
\[
    R_{p,q}(s,t)=(sp,tq),
    \qquad 0\le s,t\le1.
\]
Its pulled-back orientation form is \(pq\,ds\wedge dt\), so orientation
reverses automatically when exactly one side is negative.  Direct integration
still gives
\[
    \int_{R_{p,q}}e^{q-M}\,dA\wedge dM=p(e^q-1).
\]
Thus no separate ``unsigned area'' convention is needed.

\subsection{A closed charge path with nonzero evaluation defect}

Consider the chronological loop
\[
    \ell=(\mathsf A_p,\mathsf M_q,
           \mathsf A_{-p},\mathsf M_{-q}).
\]
Its total ACS charge is \((0,0)\), the same as the empty history
\(\varnothing\), but
\begin{align*}
    \nu_x(\ell)
      &=e^{-q}\bigl(e^q(x+p)-p\bigr)\\
      &=x+p(1-e^{-q}).
\end{align*}
Hence
\[
    \tau(\ell,\varnothing)=p(1-e^{-q}).
\]
Here \(M_*=0\), so \(\eta_*=e^{-M}dA\).  The ACS path of \(\ell\) is the
oriented rectangle boundary itself, and
\[
    \int_{C_\ell}e^{-M}\,dA
      =p-p e^{-q}
      =p(1-e^{-q}).
\]
This example simultaneously checks the closed-loop sign and shows why zero
total charge is weaker than operator neutrality.

For a self-intersecting charge path, the same computation is performed on a
singular two-chain with integer multiplicities.  Decomposing a filling into
oriented rectangles adds their signed weighted integrals; overlap is counted
with the sum of multiplicities.  This is the precise replacement for speaking
of ``the region enclosed'' by a self-intersecting loop.

\subsection{Finite contact compositions}

Write \(p=\mu h\) and \(q=\lambda k\).  On the assignment coordinate, the two
horizontal flows act by
\[
    A_h(a)=a+p,
    \qquad
    M_k(a)=e^q a.
\]
The open comparison is
\begin{align*}
    (M_k\circ A_h)(a)&=e^q(a+p),\\
    (A_h\circ M_k)(a)&=e^q a+p,
\end{align*}
and therefore
\[
    (M_k\circ A_h)(a)-(A_h\circ M_k)(a)
      =p(e^q-1).
\]
For the closed chronological word \(A_h,M_k,A_{-h},M_{-k}\), the induced
composition is
\[
    M_{-k}\circ A_{-h}\circ M_k\circ A_h,
\]
and
\begin{align*}
    (M_{-k}\circ A_{-h}\circ M_k\circ A_h)(a)
      &=e^{-q}\bigl(e^q(a+p)-p\bigr)\\
      &=a+p(1-e^{-q}).
\end{align*}
Consequently
\[
    T_{\mathrm{closed}}
      =p(1-e^{-q})
      =e^{-q}p(e^q-1)
      =e^{-q}T_{\mathrm{open}}.
\]
The equality after multiplication by \(e^q\) is a change to the common target
frame, not an equality of the two raw finite defects.

\subsection{Darboux coordinates and the Reeb field}

Assume \(\mu\lambda\ne0\) and introduce natural coordinates
\[
    \widetilde u=\mu u,
    \qquad
    \widetilde v=\lambda v.
\]
Then
\[
    \alpha=da-d\widetilde u-a\,d\widetilde v.
\]
Set
\[
    x=\widetilde v,
    \qquad y=a,
    \qquad z=a-\widetilde u.
\]
This is a global affine change of coordinates on \(\mathbb R^3\), and
\[
    dz-y\,dx
      =da-d\widetilde u-a\,d\widetilde v
      =\alpha.
\]
Thus the contact form is in the standard Darboux form.  The arithmetic
structure is not a new local contact-isomorphism class; it is the distinguished
frame \((D_u,D_v)\) and its interpretation by addition and scaling.

For completeness, the Reeb field \(R\) is determined by
\[
    \alpha(R)=1,
    \qquad
    \iota_Rd\alpha=0.
\]
Since \(d\alpha=-\lambda da\wedge dv\), its kernel is spanned by
\(\partial_u\).  Hence
\[
    R=-\frac1\mu\partial_u.
\]
In the Darboux coordinates above this is \(R=\partial_z\), as expected.

\subsection{Expanded scalar-curvature calculation}

Let \(F=F(u,v,a)\).  Expanding the two ordered derivatives gives
\begin{align*}
    D_uD_vF
       &=F_{uv}+\lambda aF_{ua}+\mu F_{av}
         +\mu\lambda F_a+\mu\lambda aF_{aa},\\
    D_vD_uF
       &=F_{vu}+\mu F_{va}+\lambda aF_{au}
         +\mu\lambda aF_{aa}.
\end{align*}
Equality of the mixed coordinate derivatives leaves
\[
    D_uD_vF-D_vD_uF=\mu\lambda F_a.
\]
Multiplying by \(du\wedge dv\) yields
\[
    \mathcal R_H(F)
       =\mu\lambda(\partial_aF)\,du\wedge dv.
\]
This is an antisymmetrized calculation on scalar fields.  It neither defines
\(\delta_H\) on one-forms nor claims that a graded differential has been
constructed.
